\begin{frame}{투영해 보기: 두 그룹이 완벽히 갈라진다}
  \textbf{5단계} --- $v_{1} = (1,1)/\sqrt{2}$에 투영:
  $p = v_{1}^{T}x = (x_{1} + x_{2})/\sqrt{2}$
  \[
    (-2,-1),\,(-1,-2) \to \frac{-3}{\sqrt{2}} \approx -2.12 \qquad
    (2,1),\,(1,2) \to \frac{+3}{\sqrt{2}} \approx +2.12
  \]
  \begin{itemize}
    \item 1차원 값 $\{-2.12, -2.12, +2.12, +2.12\}$ --- \textbf{두 그룹이
      완벽히 갈라진다}.
    \item 여기서 1차원 k-means를 $k=2$로: 중심점 $-2.12$, $+2.12$,
      할당 $[0,0,1,1]$.
    \item \kb{PCA $\to$ k-means} 파이프라인(차원 축소 $\to$ 군집화)이
      작은 예시에서 제대로 작동하는 모습.
  \end{itemize}
  \vskip0.6em
  \textbf{6단계 — 재구성 오차}: $v_{1}$만 남기고 재구성하면 각 점의 오차는
  $|x_{1}-x_{2}|/\sqrt{2} = 1/\sqrt{2} \approx 0.707$ (4개 점 모두 동일).
  평균 제곱오차 $= 4 \times 0.707^{2}/4 = 0.5$ ---
  \kq{버린 주성분의 고유값 $\lambda_{2} = 0.5$와 정확히 같음}.
  ``상위만 남길 때 잃는 정보량 = 버린 고유값의 합''을 손으로 확인.
\end{frame}
